{\large\textbf{E2: Hot Cylinder}}

\textbf{Introduction}

A uniform metal rod of length $L = 30\,\mathrm{cm}$ and radius
$r = 1\,\mathrm{cm}$ is made of an unknown metal and is kept at room temperature
$T_0 = 26.9\,\mathrm{C} = 300\,\mathrm{K}$. The metal rod is weighed to be
$m = 460\,\mathrm{g}$. Your task is to determine the thermal properties of the
unknown metal. The metal rod can be heated at one of its ends, and temperature
measurements can be performed on customizable locations along the rod. The
heater is located between $x = 0$ and $x = L_h = 3\,\mathrm{cm}$ (see fig). The
heater can be programmed by specifying a fixed power (in watts) and the
duration (in seconds) for which the heater is turned on for. Temperature
measurements are made by specifying up to five locations for the sensors along
the rod, alongside with the frequency, starting time and the ending time of the
measurements. The simulation will show the temperature readings in accelerated
”real time” (running around 10 times faster than in the real world).

\begin{center}\includegraphics[width=0.34\linewidth]{figures/EuPhO_2021_Q5_fig1.png}\end{center}

You may assume that all of the heating power goes into the rod, and that the
rod loses heat to its surroundings via heat transfer with air and black body
radiation. Heat transfer with air is linear in temperature of the rod, and can
be described by a coefficient $\alpha$ such that the heat transfer per unit area
per unit time is $\alpha(T - T_0)$. The air is well-ventilated such that
$\alpha$ can be assumed to be constant throughout the surface of the rod and
independent of the temperature of the surface. Heat loss via black body
radiation can be described using the Stefan-Boltzmann law modified with
emissivity $\beta$ such that the heat loss to radiation per unit area per unit
time is $\beta\sigma(T^4 - T_0^4)$ where
$\sigma = 5.67 \times 10^{-8}\,\mathrm{W/(m^2\,K^4)}$. Similar to $\alpha$, the
emissivity can be assumed to be constant throughout the rod, and independent of
temperature. The rod is further characterised by the thermal conductivity $k$
(such that the heat flux density along $x$ is $-k\mathrm{d}T/\mathrm{d}x$) and
the specific heat capacity $c$.

\textbf{Task}

The task is to determine the specific heat of the unknown metal, $c$ (units
$\mathrm{J/(K\,kg)}$), the thermal conductivity $k$ (units
$\mathrm{W/(m\,K)}$), and the heat loss coefficients $\alpha$ (units
$\mathrm{W/(m^2\,K)}$), and $\beta$ (dimensionless). You should aim to find the
values within $10\,\%$ of the true value. This is because there are various
sources of errors, such as Gaussian fluctuations in both defining the locations
of the sensors, and the taking the temperature measurements. The sizes of the
errors can be found by observing the fluctuations in the output.

As with all experiments, you must provide clearly labelled tables of data,
clearly labelled graphs, and sufficient formulae derivations to make it clear
what you have measured, and how you are deriving your results.

\textbf{Program interface}

Running the simulation program, named \textbf{rod}, allows performing multiple
experiments on the rod. The program will ask a sequence of prompts regarding the
setup of the experiment. For each prompt, the corresponding value(s) should be
entered, followed by pressing \textbf{return} to go the next prompt. The prompts
are as follows:

1. The heating power of the heater:\\
\texttt{Enter P (W), between 0 and 300:}

2. The duration after the start of the experiment for which the heater is
turned on for (after this time, the heater will be turned off):\\
\texttt{Enter heating duration (s), between 0 and 3600s:}

3. The starting and finishing times (after the start of the experiment) for the
temperature measurements made on the rod:\\
\texttt{Enter the starting and finishing time for the measurements (s),
separated by a space. Must be between 0 and 3600s:}

4. The time interval between two consecutive measurements that are made with
the temperature sensors:\\
\texttt{Enter dt (s), between 5 and 3600s and a multiple of 5s:}

5. The locations of the temperature sensors along the rod. The coordinates are
specified with respect to the end with the heater:\\
\texttt{Enter up to 5 locations for the sensors (in cm), between L=0 and
L=30cm, separated by spaces:}\\
Note that not entering any numbers simply means not taking any measurements.

6. The output file name for the temperature readings. Note that all of the
saved readings will also be displayed on the screen:\\
\texttt{Enter the output file name:}\\
You are advised to only use Latin letters and numbers for the name. Other
characters may or may not be allowed in the filename and in case of an invalid
filename, the readings will not be saved. The readings will be saved in a
\texttt{.txt} file with the given name in the same folder as the program.

If you enter an invalid input, a clarifying error message will be sent, and an
another opportunity for entering the input will be given.

The program will then prompt to press \textbf{return} to start the experiment,
or typing \texttt{restart} and pressing \textbf{return} to re-enter all the
experimental parameters. After continuing with the simulation, the program will
display a summary of the experimental setup, and then start printing out the
time elapsed since the heater was turned on (\texttt{t(s)}), and all the sensor
readings in the same order they were entered in the prompt (\texttt{Ti(C)},
where i corresponds to the i-th sensor).

After the simulation ends, a new experiment can be started by typing
\texttt{restart} and pressing \textbf{return}.
